\documentclass[12pt]{article}
\parindent 0.0in
\parskip 0.1in
\pagestyle{empty}
\def\R{\mbox{I}\negthinspace \mbox{R}}
\topmargin -.5in
\textheight 8.0in
\textwidth 6.0in
\oddsidemargin 0.2in
\usepackage{amsmath}
\usepackage{hyperref}

\begin{document}

\newcommand{\bm}[1]{\mbox{\boldmath$#1$}}

\begin{center}{\bf Homework Set \#5 -- Due 11:59pm, Friday, May 22.} \\ {\sc MATH 447 Spring 2026}\end{center}
\textbf{Note}: 
\begin{itemize}
\item This homework is based on Chapter 10 Part 1 \& Part 2. 
\item For your convenience, all the data files and the .tex file for this problem are included.
\item Your need to submit two files: A PDF file which contains all the answers and an R script file (with an ``.R'' extension) containing all the code.  Any answers presented in the R script file only will not be graded.
\item Your answers must be presented in the same order as the problems given in this assignment, including all graphs and {\tt R} outputs, if applicable (i.e., the graphs and outputs must be in the same order and must not be relegated to the back of your assignment). 
\item Use your judgment to include an appropriate amount of {\tt R} code and output in your answers as necessary. For example, it is not necessary for the instructor to see every single observation of the dataset you used.
\item All the {\tt R} code must be well documented and submitted as a single R script file (which has the ``.R'' extension) to Canvas. 
\item If you need an extension, you need to ask for permission before 5pm on the day the homework is due.
\end{itemize}

The grading scheme is as follows:
\begin{itemize}
\item This homework assignment will be graded out of $30$ points.
\item Each problem will be graded out of $10$ points.
\item If the R code in the .R file is missing or incomplete, at most $5$ points will be deducted.
\end{itemize}

\textbf{Textbook Problems}: Refer to \textit{An Introduction to Statistical Learning with Applications in R}, First Edition, by James, Witten, Hastie, and Tibshirani.

\begin{enumerate}
\item Principal component analysis (PCA). Use HW5code.R to answer these questions.
\begin{itemize}
\item[(a)] Present a biplot and provide a reasonable interpretation for $\phi_1$ and $\phi_2$ (the first and second principal component loading vector, respectively). It may be helpful to examine the actual values of the loading vectors to understand what they represent.
\item[(b)] By referring to both the first and second principal component, discuss important characteristics of Observation \#6083.
\item[(c)] Present a scree plot and plot of the cumulative proportion of variance explained, and decide how many principal components should be used. Be sure to justify your choice.
\end{itemize}

\item PCA and $K$-means clustering. When answering these questions, modify HW5code.R appropriately.
\begin{itemize}
\item[(a)] Interpret $\phi_1$, $\phi_2$, and $\phi_3$ (the first, second, and third principal component loading vector, respectively). It may be helpful to examine the actual values of the loading vectors to understand what they represent.
\item[(b)] Perform $K$-means clustering on $(Z1, Z2, Z3)$ using $K = 2$ and $K = 3$. Then, plot the results in three scatterplots $(Z_1, Z_2)$, $(Z_1, Z_3)$, and $(Z_2, Z_3)$ separately, and present them in your answer. Each cluster should be represented using different colors. Moreover, present the $K$ clusters as a 3D scatterplot as well. In total, you should present $3 \times 2 + 2= 8$ plots. Then, justify whether $K = 2$ or $K = 3$ is more reasonable.
\item[(c)] Using $K = 3$, interpret each cluster.
\end{itemize}

\item \# 10.7.3. For (a), use {\tt plot()}. Skip (b). Use {\tt label <- c(2,1,1,2,1,2)} for (c). \\

\end{enumerate}
\end{document}